
envoyant @@M0@@ sur @@M1@@, et @@M2@@ sur @@M3@@. Il est immédiat de voir sur (5.5.1) et (5.3.1) que c'est un @@M4@@-homomorphisme, et @@M5@@ étant intègre dans le cas considéré, on vérifie aussitôt qu'il est injectif; comme @@M6@@ est un @@M7@@-anneau, on en déduit le résultat.

Dans le @@M8@@-anneau @@M9@@, on peut définir des opérations @@M10@@ (3.2). Pour @@M11@@, la valeur de @@M12@@ calculée sur l'élément @@M13@@ de @@M14@@ est donnée par :
@@M15@@
c'est donc un élément de @@M16@@; il est défini par la relation :
@@M17@@
et les polynômes universels qui en résultent.


5.6. Soit @@M18@@ le @@M19@@-anneau engendré sur @@M20@@ par les générateurs @@M21@@, soumis aux relations @@M22@@ pour @@M23@@, @@M24@@ pour @@M25@@. On se propose de calculer les @@M26@@ en fonction des @@M27@@.

Par définition :
@@M28@@
soit :
@@M29@@
